\Page{19}{\textenglish{LAST LOOK}}{ארבע החלטות שהשתנו}{55}{63}

\vspace{10mm}

{\fontsize{17.3}{27}\selectfont

לפעמים השאלה המעניינת בקולקציה היא איזה רעיון קיבל הזדמנות להגיע לדוגמה. בעידן שבו אפשר לראות ולהשוות אפשרויות מוקדם יותר, אפשר גם לתת לפריט הפחות ברור סיבוב נוסף. במשרד הקניינות, השאלה שונה. צבע יכול להצדיק גרסה אחת של מוצר מוכר גם בלי להכתיב את כל העונה.

בפרסום, העבודה זזה מן המקום והזמן של יום הצילום אל ההחלטה אילו חלקים מהתמונה יש להפיק בכל דרך. דמות מזוהה, סרט שהמצלמה נעה בתוכו וחשבון שצריך להמשיך לפרסם דורשים שלוש עבודות שונות. בתור קונות אנחנו פוגשות שינוי אחר. אפשר לראות מראה אפשרי עוד לפני שהבגד הגיע, ועדיין להחליט מתי מראה לבדו מספיק.

\par}

\vspace{8mm}

{\color{Silver}\hrule height .55pt}

\vspace{6mm}

{\fontsize{17.3}{27}\selectfont

אלה ארבע תמורות שונות בתעשייה אחת. יהיה מעניין לראות מה מותגים יעשו עם האפשרויות שנפתחו. האם ישתמשו בהן כדי לבדוק רעיונות שהיו יקרים מדי לנסות, או יעבדו מהר יותר על מה שכבר הצליח? אופנה טומנת בחובה בחירה בין מוכר לחדש. היכולת לראות יותר מאפשרת לבחור מחדש מה ראוי לצאת לעולם.

\par}
